\chapter{Introduction}\label{sec:intro}
Intensive broadband reverberation mapping needs dense, long and homogeneous light curves, which is what the LCO
key projects in Table~\ref{tab:proposals} have delivered since 2014: daily or sub-daily multi-band imaging on the
1-m network, complemented by 2-m imaging and FLOYDS spectroscopy. Reducing this data set systematically ---
re-extracting the FLOYDS spectra with the improved pipeline, building homogeneous light curves --- starts with
knowing exactly what was observed, when, with which instrument and calibrations, and under which conditions.

This report describes a metadata database of the whole data set built from the LCO archive API
(Chapter~\ref{sec:db}), the telescope time the projects have used and its value (Chapter~\ref{sec:usage}), the
photometric and spectroscopic holdings (Chapters~\ref{sec:phot} and \ref{sec:spec}), and the cross-match of the
observed targets with the provisional AVA catalogue of AGN (Chapter~\ref{sec:sources}).

\begin{table}[h]
\centering\small
\caption{The programmes included. KEY2014A-002 and KEY2018B-001 are led by other PIs; all others are led by the
author. Filler programmes (suffix f) use the same targets.}\label{tab:proposals}
\begin{tabular}{llll}
\toprule
Proposal & Role & Semesters & Title \\
\midrule
\inputtable{tables/tab_proposals.tex}
\bottomrule
\end{tabular}
\end{table}
